% constants-source: paper/neurips2026/supplement.tex
\setcounter{figure}{0}
\renewcommand{\thefigure}{S\arabic{figure}}

\section{Data Sources and Trial Reconstruction}
\label{app:data}

\subsection{Version-pinned sources}

OpenAffect-EEG uses fixed OpenNeuro snapshots rather than a floating latest
version: MusicEEG ds002721 version 1.0.3, DENS ds003751 version 1.0.6,
EmoEEG-MC ds005540 version 1.0.7, and external validation ds006866 version
1.0.0
\citep{openneuro_ds002721,openneuro_ds003751,openneuro_ds005540,openneuro_ds006866}.
Snapshot manifests, byte counts or annex keys, and derived trial-table hashes
record the source state. Across the \claimvalue{C39} datasets we reconstruct
\claimvalue{C41} trials from \claimvalue{C40} dataset-qualified participants;
\claimvalue{C42} trials have complete harmonizable continuous targets. The core
tier contains \claimvalue{C35} participants, \claimvalue{C36} trials,
\claimvalue{C37} qualified stimulus identities, and \claimvalue{C38} supervised
trials.

\begin{table}[t]
  \caption{Source and redistribution boundaries. License entries describe the
  fixed source metadata; companion media and code retain separate terms.}
  \label{tab:licenses}
  \centering
  \small
  \setlength{\tabcolsep}{4pt}
  \begin{tabularx}{\linewidth}{@{}lXXX@{}}
    \toprule
    Dataset & OpenNeuro snapshot & Companion boundary & Benchmark release \\
    \midrule
    MusicEEG & CC0 & OSF metadata/audio: CC BY 4.0 & Metadata, splits, scores \\
    DENS & CC0 field; ODbL/DbCL conflict & Videos not mirrored & Conservative attribution \\
    EmoEEG-MC & CC0 & Stimuli/code/checkpoints separate & No EEG, DE, media, weights \\
    ds006866 & CC0 & OSF table: CC BY 4.0; source images separate & No EEG or images \\
    \bottomrule
  \end{tabularx}
  % evidence: fixed dataset_description.json files; paper/generated/table_benchmark_scope.csv
\end{table}

\begin{table}[t]
  \caption{Observed-data and probe-coverage disclosure. EEG feature counts are
  applied after a frozen trial assignment; missing signals never redefine a
  split. Probe coverage is the test fraction whose class was observed in fit.}
  \label{tab:coverage}
  \centering
  \small
  \setlength{\tabcolsep}{4pt}
  \begin{tabularx}{\linewidth}{@{}lrrXr@{}}
    \toprule
    Dataset & Labeled & EEG features & Main role & Stimulus probe \\
    \midrule
    MusicEEG & 1,240 & 1,240 & Transparent spectral audit & 0.8708 \\
    DENS & 365 & 162 & Transparent, missingness-limited & 0.9826 \\
    EmoEEG-MC & 1,260 & 1,260 & DE and foundation case study & 1.0000 \\
    ds006866 & 35,504 & 35,504 & Condition-level external audit & N/A \\
    \bottomrule
  \end{tabularx}
  % evidence: paper/generated/table_benchmark_scope.csv; paper/generated/table_identity_probes.csv
\end{table}

The DENS gap originates from unreadable public signal windows rather than an
explicit target filter, but the resulting complete-case set is not ignorable.
Only \claimvalue{C75} of \claimvalue{C76} label-eligible trials are available
(missing fraction \claimvalue{C74}); \claimvalue{C78} participants have no
available trial, and the available-versus-missing valence standardized mean
difference is \claimvalue{C77}. Missingness summaries and probe seen-class
fractions are published per dataset and seed. Subject probes in the LaBraM case study have complete test
coverage; permutation labels are shuffled only in the training partition while
the assignment, features, estimator, and test labels remain fixed.

\paragraph{MusicEEG.}
Participants listened to film-music excerpts and completed eight self-report
items \citep{daly2020musiceeg}. The benchmark retains 1,240 task trials. Public
events do not directly provide canonical audio filenames, so
\texttt{stimulus\_uid} is aligned from event codes to the public OSF Set 1 after
checking uniqueness, code range, and the availability of audio and normative
ratings for every candidate. Because the README, event description, and early
paper do not fully agree on the material set, this mapping is labeled
\texttt{audited\_inference}, not \texttt{source\_verified}. The OpenNeuro
snapshot declares CC0 and the companion OSF resource declares CC BY 4.0; audio
is not redistributed. % evidence: paper/generated/table_benchmark_scope.csv; cards/datasets/ds002721.md

\paragraph{DENS.}
DENS uses naturalistic videos and collects dimensional and categorical ratings
\citep{asif2023dens}. Events and behavior are joined by normalized stimulus
name, never row position. Several public \texttt{.fdt} lengths disagree with
their \texttt{.set} headers; unreadable full windows are recorded through an
availability mask rather than padded, wrapped, or treated as valid EEG. Public
events contain twenty local stimulus names, but the snapshot lacks the
documented stimulus and code directories, and the cited AFDI entry cannot be
reliably linked to those local identities. They are therefore valid grouping
keys but not claims of recovered source films. Conflicting CC0 and ODbL/DbCL
statements are handled under the more conservative attribution and share-alike
boundary. % evidence: paper/generated/table_benchmark_scope.csv; cards/datasets/ds003751.md

\paragraph{EmoEEG-MC.}
The dataset combines video viewing and text-guided imagery
\citep{xu2025emoeeg}. A published behavioral \texttt{video\_name} field repeats
one erroneous value for some participants and cannot be joined by row order.
The supervised subset is restricted to globally verifiable stimulus codes, and
context direction is cross-checked against source EDF trigger types. Other
reconstructed trials remain in provenance with unavailable targets. Public DE
arrays, reordered EEG, model features, and checkpoints are excluded from the
anonymous package. % evidence: cards/datasets/ds005540.md

\paragraph{External emotion regulation.}
ds006866 studies social and nonsocial images under neutral viewing, negative
viewing, and reappraisal conditions \citep{maka2026discrepancy}. It contributes
\claimvalue{C43} participants and \claimvalue{C44} image-presentation trials,
of which \claimvalue{C33} have complete valence and arousal. Public events
recover condition but not trial-level image identity, and EEGLAB annotations do
not restore it. Stimulus and joint holdout are consequently blocked by source
metadata. The OpenNeuro snapshot declares CC0; companion metadata declare CC BY
4.0, while source images can have separate terms. Neither images nor a
substitute media index are distributed. % evidence: cards/datasets/ds006866.md

\section{Target Harmonization}

MusicEEG pleasantness and energetic-arousal reports define valence and arousal;
the other datasets use their released valence and arousal fields. MusicEEG,
DENS, and ds006866 map an original one-to-nine scale by $(x-1)/8$; EmoEEG-MC
maps zero-to-seven by $x/7$. Original values, scale labels, and transformation
expressions remain in trial metadata or result provenance.
% constants-source: src/openaffect_eeg/baselines.py

The supervised target is the participant's trial-level report. Nominal stimulus
category, normative rating, and experimental condition can be inputs,
stratification variables, or audit labels, but do not replace the participant
target. Trials missing either continuous dimension remain in metadata and are
excluded only from metrics requiring both dimensions.

\section{Split Protocols}
\label{app:splits}

\paragraph{Single-factor protocols and the optimistic control.}
Complete trials are assigned before windowing. Random-trial, subject, and
stimulus holdout use 70/10/20 train/validation/test proportions within each
dataset, with at least one identity reserved for validation and test in small
groups. Random-trial splitting isolates only \texttt{trial\_uid} and records
subject and stimulus overlaps. Strict single-factor splits require an empty
train--test intersection for the held identity.
% constants-source: src/openaffect_eeg/splits.py; scripts/build_strict_splits.py

\paragraph{Joint holdout.}
Subjects and stimuli are independently assigned 60/20/20. A trial is retained
only when both identities receive the same partition; cross-partition edges are
marked \texttt{excluded}. This sacrifices coverage in sparse subject--stimulus
graphs to obtain a genuine double cold start. The ratio was fixed before model
evaluation because smaller validation identity sets did not provide adequate
validation coverage across seeds.
% constants-source: scripts/build_strict_splits.py

\paragraph{Support-matched identity-overlap control.}
For each joint-holdout seed, the control freezes the exact test trial IDs and
matches the joint split's train and validation row counts. It then samples those
training-side rows so that both test subjects and test stimuli occur in the fit
partition. Consequently, the control preserves test support and sample counts
while changing training-side identity overlap and the associated graph. It is
an estimand for the practical effect of available identity structure, not a
claim that all higher-order graph properties are identical.
Across seeds, test sets contain \claimvalue{C79}--\claimvalue{C80} trials from
\claimvalue{C81}--\claimvalue{C82} participants and exactly
\claimvalue{C83} held-out stimuli.
% constants-source: src/openaffect_eeg/reviewer_controls.py; scripts/build_support_matched_control.py

% constants-source: paper/figure_plan_zh.json
\placeholderfigure{figures/figureS_protocol_audit.pdf}{2.20in}{Protocol-level
audit retained as supporting evidence. \textbf{A:} five fixed split seeds and
their mean macro-CCC for semantic prior, direct EEG, and prior-plus-EEG residual;
random-trial splitting is an optimistic control. \textbf{B:} pretrained versus
fixed-random LaBraM comparisons and support-matched identity-overlap effects.
These estimates motivate the response-surface analysis but do not by themselves
separate participant calibration from repeated-stimulus exposure.}{fig:protocol-audit}

\paragraph{Context and dataset holdout.}
Directional EmoEEG-MC transfer first holds out subjects. Only source-context
trials enter training and validation, and only target-context trials from unseen
subjects enter test. Video-to-imagery and imagery-to-video are evaluated
separately. Context and dataset holdout treat a complete domain as a group.
% constants-source: src/openaffect_eeg/splits.py; scripts/build_context_transfer_splits.py

All main representation audits use \claimvalue{C47} fixed seeds. Each assignment
stores trial roles and an audit JSON. Builders check UID uniqueness, nonempty
identity fields, legal roles, and the required empty identity intersections.
Models apply target and feature-availability masks only after loading an
assignment, so representation missingness cannot silently alter the split.

\subsection{Held-out deployment-selection utility}
\label{app:selection-utility}

To test a practical consequence rather than assume that ranking changes help,
we compare a validation-cell selector with one global validation-mean selector.
Both receive identical predictions from seven candidates: a personalized
EEG-free prior and direct/combined predictions from spectral features and frozen
pretrained/random LaBraM. Candidates, dose grids, role seed and rules are fixed
before this analysis. No candidate is removed based on test performance.
% constants-source: configs/deployment_utility_v5.json

Participant and stimulus identities have disjoint fitting, inner tuning,
selector-validation and test roles. All trials containing the other phase's
participant or stimulus identities are quarantined, including resource trials.
Assignments use constant placeholders for target-block labels. Ridge tuning
uses only inner tuning outcomes; selector choices are hashed before test
evaluation. Five resource seeds share the same role partition and target
support, not five independent cohorts. The EEG-free and EEG predictors receive
identical calibration labels and stimulus-label resources.
% constants-source: configs/deployment_utility_v5.json; src/openaffect_eeg/deployment_benchmark.py

EmoEEG-MC has 48 trials from six participants and eight stimuli in each of the
two phases. MusicEEG has 25 validation trials (six participants, eight stimuli)
and 23 test trials (six participants, nine stimuli), after maximum-dose
feasibility filtering. The sparse support limits precision and coverage. All
2,000 paired crossed draws are valid. They condition on fits and validation
choices, not selector-training uncertainty. Source, assignment and prediction
hashes are verified, and a separate centered-moment implementation reproduces
the scores and intervals. This is developer verification, not an independent
human reproduction.
% evidence: paper/generated/deployment_utility_v5/ds005540_verification.json; paper/generated/deployment_utility_v5/ds002721_verification.json
% constants-source: configs/deployment_utility_v5.json

\begin{table}[t]
\centering\small
\caption{Held-out deployment-selection comparison. Both policies use the same validation predictions. Differences are conditioned minus global, uniformly averaged across resource cells; intervals are paired crossed 95\% intervals conditional on fits and choices. CCC is primary; MAE is secondary and lower is better.}
\label{tab:selection-utility}
\begin{tabular}{@{}llrrl@{}}
\toprule
Dataset & Metric & Global & Conditioned & Difference [interval] \\
\midrule
EmoEEG-MC & CCC & 0.34221 & 0.34477 & +0.00256 [-0.00113, 0.00601] \\
EmoEEG-MC & MAE & 0.19523 & 0.19464 & -0.00059 [-0.00103, -0.00013] \\
MusicEEG & CCC & 0.24712 & 0.24600 & -0.00113 [-0.01101, 0.01027] \\
MusicEEG & MAE & 0.29405 & 0.29021 & -0.00384 [-0.00675, -0.00094] \\
\bottomrule
\end{tabular}
% evidence: paper/generated/deployment_utility_v5/ds005540.json; paper/generated/deployment_utility_v5/ds002721.json
\end{table}


Neither primary CCC interval excludes zero. Both secondary MAE differences
are small and negative, but no practical-importance threshold was established;
they cannot replace the primary endpoint or identify EEG added value. The
global choices are random-LaBraM combined prediction for EmoEEG-MC and spectral
combined prediction for MusicEEG. These choices describe this candidate set and
validation block, not population-level model superiority. This analysis uses
datasets already examined during project development: role isolation protects
the current selection procedure, but is not untouched external confirmation.
% evidence: paper/generated/deployment_utility_v5/ds005540.json; paper/generated/deployment_utility_v5/ds002721.json

\subsection{Identity-exposure response surface}
\label{app:exposure-surface}

Identity availability is not represented as one binary variable. A deployment
state can separately specify whether the participant ID is known, whether
unlabelled target-participant data are available, how many labelled calibration
trials are available, whether calibration comes from the same session, and
whether the test stimulus, dataset, and device occurred during fitting. The
present response surface varies exactly two labelled resources: $p$, complete
calibration trials per target participant, and $s$, other-participant ratings
per test stimulus (or experimental condition where stimulus identity is
unavailable). Released normative stimulus metadata, where present, is held fixed
and reported as a separate prior rather than counted as $s$. It fixes unlabelled
target-participant adaptation at zero. Session
and device are not separately isolated by the available metadata and are
therefore blocked claims, not implicit successes.

For each seed, we begin from one strict participant--stimulus joint assignment.
The test set is restricted once to participants and stimuli supporting the
predeclared maximum doses; its trial IDs are then hash-fixed for every $(p,s)$
cell. Participant calibration candidates contain a test participant under a
non-test stimulus. They never enter population-model fitting. Repeated-stimulus
candidates contain a test stimulus rated by a non-test participant. Dose samples
are deterministic and nested, so every lower dose is a subset of every higher
dose on the same axis.

Adding $s$ repeated-stimulus rows would otherwise increase sample size. Each
added row therefore replaces one base-training donor selected deterministically
within context and nearest in the two training targets. Population-training size
is exactly constant across cells, while validation and test trials are unchanged.
Donor selection reads training labels only. The intervention intentionally
changes which stimulus identities are represented during fitting; it does not
claim all graph or target-distribution properties are causally matched.

% constants-source: src/openaffect_eeg/identity_exposure.py; configs/identity_exposure_response_v1.yaml
Let $f_{s}$ be a population predictor fit with stimulus exposure dose $s$. For a
target participant $i$, labelled calibration produces only the post-fit offset
\begin{equation}
  \widehat{\boldsymbol\delta}_{i,p}
  = \frac{1}{p}\sum_{j\in\mathcal C_{i,p}}
  \left(\mathbf y_{ij}-f_s(\mathbf x_{ij})\right),
\end{equation}
with zero offset at $p=0$. Test prediction is
$f_s(\mathbf x_{i,t})+\widehat{\boldsymbol\delta}_{i,p}$. Thus a nonzero $p$
cell is explicitly a known-participant, labelled-calibration result; it is not
subject independent. The same construction is applied to the prior-plus-EEG
predictor and to the EEG-free prior with exactly the same calibration rows.
Population and personalized outputs are stored separately.

\begin{table}[t]
\centering\small
\setlength{\tabcolsep}{3pt}
\caption{Maximum-resource predictions with resource-matched controls. P: prior; PC: calibrated prior without EEG; E: prior plus population EEG residual; EC: calibrated E. Brackets are paired crossed-identity 95\% pointwise intervals. These are not equivalence tests.}
\label{tab:paired-added-value}
\begin{tabular}{@{}lrrrrl@{}}
\toprule
Representation & P & PC & E & EC & EC $-$ PC [95\% CI] \\
\midrule
\multicolumn{6}{l}{\textit{EmoEEG-MC (p=8, s=8)}} \\
DE & 0.4809 & 0.5863 & 0.4416 & 0.5599 & -0.0265 [-0.0640, 0.0178] \\
EEGNet & 0.4809 & 0.5863 & 0.4863 & 0.5955 & 0.0091 [-0.0059, 0.0307] \\
LaBraM pretrained & 0.4809 & 0.5863 & 0.4676 & 0.5748 & -0.0115 [-0.0496, 0.0298] \\
LaBraM random & 0.4809 & 0.5863 & 0.4629 & 0.5711 & -0.0152 [-0.0553, 0.0179] \\
CBraMod pretrained & 0.4809 & 0.5863 & 0.4590 & 0.5686 & -0.0177 [-0.0472, 0.0235] \\
CBraMod random & 0.4809 & 0.5863 & 0.4817 & 0.5840 & -0.0023 [-0.0327, 0.0262] \\
EMOD frozen & 0.4809 & 0.5863 & 0.4738 & 0.5819 & -0.0045 [-0.0359, 0.0258] \\
EMOD random & 0.4809 & 0.5863 & 0.4530 & 0.5687 & -0.0176 [-0.0669, 0.0278] \\
\midrule
\multicolumn{6}{l}{\textit{MusicEEG (p=8, s=4)}} \\
Band power & 0.1740 & 0.2930 & 0.1799 & 0.2988 & 0.0058 [-0.0131, 0.0253] \\
LaBraM pretrained & 0.1740 & 0.2930 & 0.2085 & 0.3070 & 0.0141 [-0.0332, 0.0569] \\
LaBraM random & 0.1740 & 0.2930 & 0.1805 & 0.2941 & 0.0011 [-0.0531, 0.0587] \\
\bottomrule
\end{tabular}
\end{table}


\paragraph{Matched model extension and reproducibility.}
The EmoEEG-MC extension evaluates EEGNet and frozen EMOD on every existing
assignment, asserting identical test IDs, targets, priors, and calibrated priors.
EEGNet selects its epoch count by validation MAE, then refits on train plus
validation; target-participant calibration does not update network weights.
Direct and residual models are fitted separately. Population fits are reused
across participant doses only after verifying unchanged fit assignments.
The compact model uses the previously specified custom EEGNet-style architecture with zero
adversarial loss. All preprocessing statistics are estimated from its fit data.

% constants-source: scripts/extract_emod_audit_features.py
EMOD uses official repository revision
\texttt{3a219f51ba4f90b9e2d22824cfb90dde77a9875c}.
Its checkpoint has legacy pooling/projection names; the adapter records a
one-to-one name mapping and then requires a strict match of all tensors.
It uses the benchmark's 100-Hz raw archive, three nonoverlapping ten-second
crops, and mean pooled official projection embeddings. Channels absent from the
official channel vocabulary are excluded and listed in the extraction manifest.
Input scaling/preprocessing and head adaptation differ from the paper's full
fine-tuning pipeline. This audits the supplied frozen weights, not a claim that
EMOD's original reported SOTA performance was reproduced. Neither third-party
weights nor raw feature arrays are redistributed.

\paragraph{Failed native-task learning check.}
A separate bounded SEED-V check used the official SEEDV model branch and
checkpoint backbone with 62 channels, 200 Hz, microvolt inputs and one-second
windows. Within each participant/session, the first five, middle five and last
five trials formed train, validation and test. Across 48 source files there
were 117,744 windows; 40,352 test windows represented 240 trials from 16
participants. A single fixed seed, 100 epochs, and validation-kappa checkpoint
selection yielded test balanced accuracy 0.20 and kappa 0: all test windows
received one class. Prepared-input hashes, test identities, finite gradients
and recomputed metrics were verified, but native-task learning was not
recovered. The released SEEDV branch uses a large multilayer head, whereas the
paper describes a linear classifier; the per-task selected fine-tuning recipe
is not fully specified. This failed compatibility check neither reproduces nor
refutes the published scores. It cannot validate the frozen regression audit
as a full SOTA test. CUDA pooling backward is nondeterministic despite fixed
seeds; this is recorded with the aggregate training history.
% evidence: paper/generated/competitive_revision_v3/emod_native/native_task_check.json; paper/generated/competitive_revision_v3/emod_native/training_history.json
% constants-source: scripts/run_emod_seedv_check_v3.py; scripts/prepare_emod_seedv_v3.py

\paragraph{Training-only diagnosis and repair.}
We retained that failed run and reproduced its numerical failure on 50 fixed
training windows, ten per class, without validation/test-based parameter
selection. With classifier learning rate $10^{-3}$, the first update increased
the first dense layer's output standard deviation from 0.773 to 32.280;
by step 20, the second hidden ELU produced constant outputs. The official SEEDV
head flattens 39,680 features into a 2,480-unit layer. Keeping the architecture,
inputs, initialization, optimizer and backbone learning rate $5\times10^{-4}$
unchanged, reducing only the classifier learning rate to $10^{-5}$ prevented
saturation and fitted the diagnostic batch by step 50. This establishes a
configuration-specific numerical remedy, not an emotion-mechanism finding.
All 48 training-label files matched the official constants. Reprocessing 100
training windows across five source files, including the format-fallback and
largest-amplitude cases, matched the prepared float32 inputs exactly; this is
selected-file parity, not exhaustive independent data validation.
% evidence: paper/generated/emod_repair_v4/native_task_check.json; paper/generated/emod_repair_v4/original_lr_training_trace.json; paper/generated/emod_repair_v4/repaired_lr_training_trace.json
% constants-source: configs/emod_native_repair_v4.json; scripts/check_emod_preparation_parity.py

The repaired run completed the fixed 100-epoch budget. Validation kappa selected
epoch 21, after which the test set was evaluated once. Recomputed test BACC was
38.62\%, kappa 0.2348 and weighted F1 39.43\%; all five classes were predicted.
The published five-run BACC is $41.24\pm0.86\%$ \citep{chen2026emod}, leaving
a 2.62 percentage-point gap to that mean. This single-seed check demonstrates
recovered learning, not statistical equivalence or a five-run reproduction.
It uses within-participant classification, not the self-report exposure
response surface, and cannot establish unseen-user transfer or pretrained
benefit without a matched random-initialization comparison. Aggregate traces,
input/source hashes and a fail-fast input-sensitivity check accompany the
repair. The original failed results remain archived.
% evidence: paper/generated/emod_repair_v4/native_task_check.json; paper/generated/emod_repair_v4/training_history.json
% external-values: chen2026emod, Table 2

\paragraph{Standard implementation and stronger calibration control.}
The added standard regression model uses Braindecode 1.2.0 EEGNet with
$F_1=8$, $D=2$, $F_2=16$, temporal kernels 64/16, dropout 0.25, and two raw
regression outputs. AdamW uses learning rate 0.001 and weight decay 0.0001;
validation MAE selects among at most 100 epochs with patience 12. All 25
population states and 125 dose cells use the archived assignments. Its 50
direct/residual fits have finite logged objectives and gradients; selected
epochs range from 1 to 24, with 14 fits selecting epoch one. Synthetic learning
and training-only tiny-subset fitting checks pass, but do not demonstrate
optimal benchmark tuning or held-out emotion information.
% constants-source: configs/competitive_revision_v3.json
% evidence: paper/generated/competitive_revision_v3/standard_training_diagnostics/training_verification.json; paper/generated/competitive_revision_v3/standard_training_diagnostics/learning_checks.json

An additional empirical calibration control estimates within-participant
residual variance $\sigma^2$ and nonnegative noise-corrected between-participant
mean variance $\tau^2$ from population-fit data. It multiplies every calibrated
offset by $w_p=p\tau^2/(p\tau^2+\sigma^2)$, with $w_0=0$. These are heuristic
prediction weights, not identified latent variance components. The same weights
apply to $PC$ and $EC$. Table~\ref{tab:v3-calibration} reports uniform cell means;
the historical dose-weighted integrals remain separate secondary summaries.
Neither calibration rule establishes a positive surface-average EEG increment
for the displayed models at the reported precision.
% evidence: paper/generated/competitive_revision_v3/table_provenance.json

\begin{table}[t]
\centering\small
\caption{Implementation and calibration sensitivity. Entries are uniform-grid mean matched EEG added value, $\mathrm{CCC}(EC)-\mathrm{CCC}(PC)$, with 95\% paired crossed intervals conditional on fitted predictions. Shrinkage weights are estimated only from population-fit residuals and shared by EEG and non-EEG offsets. These are not maximum-dose scores or equivalence tests.}
\label{tab:v3-calibration}
\begin{tabular}{@{}llcc@{}}
\toprule
Dataset & Representation & Mean offset & Shrunk offset \\
\midrule
EmoEEG-MC & Compact EEGNet-style & 0.0089 [-0.0034, 0.0254] & 0.0067 [-0.0064, 0.0253] \\
EmoEEG-MC & Standard EEGNet & 0.0006 [-0.0036, 0.0054] & 0.0006 [-0.0044, 0.0060] \\
EmoEEG-MC & LaBraM & -0.0100 [-0.0392, 0.0279] & -0.0098 [-0.0419, 0.0308] \\
EmoEEG-MC & Frozen EMOD & -0.0001 [-0.0280, 0.0247] & -0.0005 [-0.0302, 0.0263] \\
MusicEEG & Band power & 0.0042 [-0.0094, 0.0188] & 0.0052 [-0.0269, 0.0335] \\
MusicEEG & LaBraM & 0.0182 [-0.0161, 0.0452] & 0.0247 [-0.0176, 0.0568] \\
\bottomrule
\end{tabular}
\end{table}

\begin{table}[t]
\centering\small
\caption{Companion error contrasts for the same matched resources. Entries are uniform-grid mean $\mathrm{MAE}(EC)-\mathrm{MAE}(PC)$ with paired crossed 95\% intervals. Negative values favor the EEG branch; positive values mean larger error. These marginal intervals do not control error across all representations and calibration variants.}
\label{tab:v3-mae}
\begin{tabular}{@{}llcc@{}}
\toprule
Dataset & Representation & Mean offset & Shrunk offset \\
\midrule
EmoEEG-MC & Compact EEGNet-style & -0.0008 [-0.0030, 0.0011] & -0.0002 [-0.0032, 0.0023] \\
EmoEEG-MC & Standard EEGNet & 0.0004 [-0.0004, 0.0012] & 0.0005 [-0.0005, 0.0015] \\
EmoEEG-MC & LaBraM & 0.0059 [0.0001, 0.0114] & 0.0056 [-0.0009, 0.0122] \\
EmoEEG-MC & Frozen EMOD & 0.0032 [-0.0025, 0.0100] & 0.0029 [-0.0029, 0.0099] \\
MusicEEG & Band power & 0.0013 [-0.0021, 0.0044] & 0.0033 [-0.0045, 0.0108] \\
MusicEEG & LaBraM & -0.0006 [-0.0071, 0.0067] & -0.0010 [-0.0100, 0.0096] \\
\bottomrule
\end{tabular}
\end{table}


\paragraph{Revised uncertainty.}
% evidence: paper/generated/identity_exposure_v2/complete/added_value_statistics.json
The paired tables use 2,000 shared crossed-identity draws across fixed splits.
The same draws compare every predictor, dose, and model; a degenerate draw is
not repaired by silently dropping a seed. The tables expose valid draw counts.
The legacy resource-effect table below preserves the original seedwise bootstrap
for traceability and is not the revised estimator used in the main text.
In particular, MusicEEG's revised joint interval crosses zero. For ds006866,
each fixed split holds out one condition; its condition factor cannot support
a within-split resampling claim about new conditions.

% evidence: paper/generated/identity_exposure_v2/complete/music_dose_support.json
MusicEEG's frozen grid is five participant doses by four stimulus doses,
ending at $s=4$. Raising $s$ to eight cannot preserve the same test support in
four of the five seeds: their available other-participant observations are
insufficient. Only one seed supports that extension without dropping test
trials. We do not fill unavailable cells or compare a changed test set as if
it were the same response surface.

% constants-source: src/openaffect_eeg/identity_exposure.py; src/openaffect_eeg/audit_cli.py
The artifact emits one deterministic claim card per cell. It records ID and
label requirements and blocks latent-emotion disentanglement, causal identity
removal, unseen-participant claims when $p>0$, unseen-stimulus claims when
$s>0$, and session/dataset/device transfer unless those factors are separately
held out. This is rule-based compilation from split lineage, not free-form
language-model judgment.

\begin{table}[t]
  \caption{Legacy fixed-support identity-exposure summary (original seedwise
  bootstrap, retained for traceability). Values are macro-CCC changes
  from the zero-resource cell; brackets are crossed participant--stimulus 95\%
  intervals. For ds006866, $S$ denotes experimental-condition exposure because
  source metadata do not expose trial-level image identity.}
  \label{tab:exposure-summary}
  \centering
  \scriptsize
  \setlength{\tabcolsep}{2.5pt}
  \resizebox{\linewidth}{!}{%
  \begin{tabular}{@{}llcccccc@{}}
    \toprule
    Dataset & Representation & Zero & Maximum & Surface AUC & $\Delta P$ & $\Delta S$ & $\Delta(P{+}S)$ \\
    \midrule
    EmoEEG-MC & LaBraM pretrained
      & \claimvalue{C105} & \claimvalue{C106} & \claimvalue{C107}
      & \claimvalue{C111} [\claimvalue{C112}, \claimvalue{C113}]
      & \claimvalue{C114} [\claimvalue{C115}, \claimvalue{C116}]
      & \claimvalue{C117} [\claimvalue{C118}, \claimvalue{C119}] \\
    MusicEEG & LaBraM pretrained
      & \claimvalue{C108} & \claimvalue{C109} & \claimvalue{C110}
      & \claimvalue{C120} [\claimvalue{C121}, \claimvalue{C122}]
      & \claimvalue{C123} [\claimvalue{C124}, \claimvalue{C125}]
      & \claimvalue{C126} [\claimvalue{C127}, \claimvalue{C128}] \\
    ds006866 & Band power
      & -- & -- & --
      & \claimvalue{C137} [\claimvalue{C138}, \claimvalue{C139}]
      & \claimvalue{C140} [\claimvalue{C141}, \claimvalue{C142}]
      & \claimvalue{C143} [\claimvalue{C144}, \claimvalue{C145}] \\
    \bottomrule
  \end{tabular}}
  % evidence: paper/generated/identity_exposure_v1/analysis/table_identity_exposure_effects.csv
  % evidence: paper/generated/identity_exposure_v1/analysis/table_identity_exposure_deployment.csv
\end{table}

% evidence: paper/generated/identity_exposure_v1/analysis/table_identity_exposure_seed_scores.csv
\placeholderfigure{figures/figureS_exposure_corners.pdf}{2.35in}{Deployment
corners for the three main representations. Large markers and connecting lines
show seed means; faint points are the five fixed split seeds. ``Participant''
uses the maximum labelled calibration dose with unseen stimuli, ``Stimulus''
uses the maximum repeated-stimulus dose with unseen participants, and ``Both''
uses both resources. The large MusicEEG spread is why endpoint means are not
treated as independent replications.}{fig:exposure-corners}

% evidence: paper/generated/identity_exposure_v1/analysis/table_identity_exposure_effects.csv
\placeholderfigure{figures/figureS_exposure_effects.pdf}{2.35in}{Maximum-dose
minus zero-resource changes. Bars distinguish transparent features, pretrained
LaBraM, and fixed-random LaBraM; whiskers are crossed participant--stimulus
bootstrap 95\% intervals. EmoEEG-MC supports positive marginal and joint
changes under the legacy estimator. The revised shared-identity interval for
MusicEEG's joint change crosses zero; this legacy plot is not evidence of a
confirmed external resource effect.}{fig:exposure-effects}

\begin{table}[t]
  \caption{Dataset--protocol applicability matrix. ``Blocked'' denotes missing
  source identity metadata, not model failure.}
  \label{tab:protocol-matrix}
  \centering
  \small
  \setlength{\tabcolsep}{4pt}
  \begin{tabular}{@{}lcccccc@{}}
    \toprule
    Dataset & Random & Subject & Stimulus & Joint & Context & Dataset \\
    \midrule
    MusicEEG & Yes & Yes & Yes & Yes & N/A & Yes \\
    DENS & Yes & Yes & Yes & Yes & N/A & Yes \\
    EmoEEG-MC & Yes & Yes & Yes & Yes & Yes & Yes \\
    ds006866 & Yes & Yes & Blocked & Blocked & Yes & N/A \\
    \bottomrule
  \end{tabular}
  % evidence: cards/evaluations/evaluation_matrix.tsv
\end{table}

\begin{table}[t]
  \caption{Executed analysis coverage. This table is intentionally distinct
  from protocol applicability: ``supported'' does not imply that every model
  was run in that cell. V/I denotes directional video--imagery transfer.}
  \label{tab:executed-matrix}
  \centering
  \small
  \begin{tabularx}{\linewidth}{@{}l l X@{}}
    \toprule
    Dataset & Representation or analysis & Executed protocols \cr
    \midrule
    EmoEEG-MC & DE, CBraMod, LaBraM affect heads & Random, subject, stimulus, joint, V/I; exposure grid \cr
    EmoEEG-MC & EEGNet & Joint; complete exposure grid \cr
    EmoEEG-MC & EMOD frozen and random & Complete exposure grid \cr
    MusicEEG & Band power, LaBraM pretrained/random & Complete feasible exposure grid \cr
    MusicEEG, DENS, EmoEEG-MC & Transparent identity probes & Dataset, subject, stimulus as applicable \cr
    ds006866 & Channel/global bandpower & Subject holdout and condition probe \cr
    \bottomrule
  \end{tabularx}
  % evidence: configs/recompute_neurips_ed_v1.yaml; paper/generated/registry/artifact_manifest.json
\end{table}

\section{EEG Representations and Models}
\label{app:models}

\paragraph{Transparent spectra.}
Core transparent features use delta, theta, alpha, beta, and gamma DE or Welch
log power. Global views aggregate channels; channel views retain the
channel-by-band layout. Standardization is fitted on training data only. Ridge
regularization is selected by validation MAE and refitted on train plus
validation. Single-band, leave-one-band, single-region, and leave-one-region
views test predictive sensitivity, not source localization or causal anatomy.

\paragraph{CBraMod.}
Official pretrained and same-architecture random encoders receive identical
EmoEEG-MC trials. Each 30-second, 64-channel, 200-Hz trial is organized into
one-second patches; channel and time tokens are mean-pooled into a 200-dimensional
trial representation. Source commit, checkpoint SHA-256, trial UIDs, and
feature hashes accompany the artifact. The encoder is frozen and only the
common downstream head is trained.
% constants-source: src/openaffect_eeg/foundation.py; cards/models/cbramod.md

\paragraph{LaBraM.}
LaBraM uses the official channel-location map and base checkpoint. Thirty
one-second patches are processed as three eight-patch windows and one six-patch
window, then pooled with patch-count weighting into a 200-dimensional trial
representation. Pretraining masks and decoder heads are excluded, while the
official final normalization path is retained. The random control instantiates
the same architecture without loading a checkpoint.
% constants-source: src/openaffect_eeg/labram.py; cards/models/labram.md

\paragraph{Historical compact EEGNet-style control.}
The custom task-specific control is trained from EmoEEG-MC waveforms resampled with
anti-aliasing to 100 Hz. Each epoch uses one deterministic 10-second training
crop; validation and test average three nonoverlapping crops. The network uses
temporal convolution, depthwise spatial convolution, separable temporal
convolution, and an eight-dimensional latent. It tests architecture dependence,
not an exhaustive architecture search.
% constants-source: scripts/run_eegnet_disentanglement.py; cards/models/eegnet.md

\begin{table}[t]
  \caption{Representation controls and training state.}
  \label{tab:representations}
  \centering
  \small
  \begin{tabularx}{\linewidth}{@{}lXXX@{}}
    \toprule
    Representation & Trial view & Pretraining control & Downstream state \\
    \midrule
    DE / Welch & Global, channel, band, region & N/A & Validation-selected Ridge \\
    CBraMod & Frozen 200-D embedding & Same architecture, random & Frozen encoder \\
    LaBraM & Frozen 200-D embedding & Same architecture, random & Frozen encoder \\
    EEGNet & Raw deterministic crops & Zero adversary & End-to-end control \\
    EMOD & Frozen 128-D embedding & Same architecture, random & Validation-selected Ridge \\
    \bottomrule
  \end{tabularx}
  % constants-source: cards/models/de_ridge.md; cards/models/cbramod.md; cards/models/labram.md; cards/models/eegnet.md
  % evidence: paper/generated/identity_exposure_v2/complete/execution_evidence.json
\end{table}

\section{Predictors, Debiasing, and Identity Probes}

\paragraph{Predictors.}
\texttt{direct\_eeg} predicts normalized targets from EEG.
\texttt{semantic\_stimulus\_prior} combines one-hot nominal category, context,
and source stimulus-intensity code with lower-cased TF--IDF description
unigrams/bigrams (at most 1,024 features, sublinear term frequency). A Ridge
model is fitted to one mean target row per training stimulus; alpha is selected
from $\{0.01,0.1,1,10,100,1000\}$ by trial-level validation MAE. The intensity
code identifies source stimulus variants and never uses a held-out participant's
rating. The historical registry identifier
\texttt{disentangled\_prior\_plus\_eeg\_residual} uses a training-only
stimulus report mean for seen stimuli and the semantic estimate for unseen
stimuli, forms leave-one-trial-out residual targets, predicts them from EEG, and
adds the hybrid prior. All branches share the assignment and validation rule;
none reads test targets. The identifier denotes arithmetic residualization and
does not claim causal or latent-state disentanglement.
% constants-source: scripts/run_disentanglement_baseline.py; src/openaffect_eeg/disentanglement.py

\begin{table}[t]
  \caption{Optimization and selection budget. Every choice is made on the
  validation partition and then reused for the paired test comparison.}
  \label{tab:model-budget}
  \centering
  \small
  \begin{tabularx}{\linewidth}{@{}lXl@{}}
    \toprule
    Component & Validation search & Training budget \cr
    \midrule
    Ridge prior/EEG/residual heads & Six alphas from $10^{-2}$ to $10^3$ & Closed-form per seed/protocol \cr
    Linear nuisance controls & Five strengths from 0 to 1 & Closed-form per strength \cr
    Gradient-reversal residual & Six adversary strengths & At most 200 epochs; patience 25 \cr
    Adversarial EEGNet & Three adversary strengths & At most 100 epochs; patience 12 \cr
    Subject-subspace audit & Eight ranks, three direction types & Head refit at every point \cr
    \bottomrule
  \end{tabularx}
  % constants-source: scripts/run_disentanglement_baseline.py; scripts/run_adversarial_disentanglement.py; scripts/run_eegnet_disentanglement.py; scripts/run_subject_subspace_intervention.py
\end{table}

\paragraph{Debiasing counterfactuals.}
We evaluate stimulus centering, conditional covariance penalties,
response-preserving stimulus-subspace projection, gradient-reversal residual
networks, and adversarial EEGNet. Every method contains a zero-strength setting;
adversarial strength, projection rank, and epoch are selected by validation
residual MAE only. Test-time residual prediction does not use stimulus ID. A
successful counterfactual must jointly improve target performance and reduce
the relevant identity effect. Variance shrinkage or inconsistent seed
direction is reported as a negative result.

% constants-source: paper/figure_plan_zh.json
\placeholderfigure{figures/figure4.pdf}{2.20in}{Subject-subspace intervention
audit. \textbf{A} plots affect macro-CCC and \textbf{B} participant balanced
accuracy across the fixed removal-rank grid. Faint trajectories are the five
evaluation seeds; opaque lines and markers are their means. Triangles remove
subject-aligned directions, squares remove equal-rank PCA directions, and
circles remove equal-rank random directions. The vertical dashed line is the
nested zero-rank control; the horizontal dotted line in \textbf{B} is chance.
The figure exposes an accuracy--identity trade-off and does not identify
participant structure as a pure nuisance factor.}{fig:debiasing-audit}

\paragraph{Identity probes.}
Dataset identity is predicted under subject holdout; stimulus identity under
subject holdout is evaluated only on classes observed in training; subject
identity under stimulus holdout is likewise evaluated on training-seen people.
Each probe uses a matched training-label permutation and reports balanced
accuracy, class count, uniform chance, coverage, and per-seed values.

\paragraph{Subject-subspace intervention.}
For each stimulus-holdout seed, between-subject centroid differences are
estimated on training data, decomposed by SVD, and the leading directions are
removed from fit and test LaBraM representations. The requested removal grid is
$\{0,1,2,4,8,16,24,29\}$. At every rank, we also remove the same number of
training-estimated PCA directions and seeded random orthonormal directions.
Semantic, direct, residual, and subject-probe heads are refitted with the same
validation-only selection. The zero-removal point is the nested control.
This intervention destroys linearly decodable subject structure but can also
remove participant-related affect variation; it is not interpreted as a pure
nuisance projector.
% constants-source: src/openaffect_eeg/reviewer_controls.py; scripts/run_subject_subspace_intervention.py

\section{Statistical Analysis}

Continuous predictions report valence, arousal, macro CCC, MAE, and RMSE.
Identity tasks report balanced accuracy and, where available, macro-F1. Model
comparisons are paired on identical trial-level predictions. In the legacy
non-surface analyses, the hierarchical
bootstrap first resamples evaluation seeds and then complete held-out identity
clusters; joint holdout separately resamples crossed subject and stimulus
factors. Those historical intervals use 2,000 resamples and a fixed statistical
seed. This is not the revised primary response-surface estimator, which holds
split seeds fixed and shares crossed identity weights across every model/cell.
% constants-source: configs/recompute_neurips_ed_v1.yaml; scripts/analyze_paired_predictions.py

With \claimvalue{C47} seeds, the exact two-sided sign-flip test has a minimum
nonzero $p$-value of \claimvalue{C45}. We therefore interpret bootstrap
intervals as effect uncertainty and the sign-flip test as a finite-seed
resolution statement. Historical targeted robustness analyses comprise the
support-matched CCC contrast, subject-subspace endpoint CCC contrast,
unseen-stimulus LaBraM pretraining contrast, and external condition-prior paired
contrasts. These were selected during benchmark development and reviewer-driven
audit rather than preregistered. Other reported cells are exploratory. Seeds
represent Monte Carlo split variability; participant, stimulus, or crossed
cluster resampling supplies the corresponding generalization uncertainty.
Intervals are pointwise and are not presented as simultaneous family-wise
coverage.

\begin{table}[t]
  \caption{Selected paired macro-CCC comparisons. Intervals are hierarchical
  paired-bootstrap 95\% intervals.}
  \label{tab:paired}
  \centering
  \small
  \begin{tabularx}{\linewidth}{@{}lXrrX@{}}
    \toprule
    Comparison & Protocol & $\Delta$CCC & Interval & Scope \\
    \midrule
    LaBraM pretrained $-$ random & Stimulus & \claimvalue{C04} &
      [\claimvalue{C05}, \claimvalue{C06}] & Known subjects, new stimuli \\
    LaBraM pretrained $-$ random & Subject & \claimvalue{C07} & Includes zero & New subjects \\
    LaBraM pretrained $-$ random & Joint & \claimvalue{C08} & Includes zero & New subjects and stimuli \\
    Identity overlap $-$ joint, direct EEG & Matched support & \claimvalue{C61} &
      [\claimvalue{C62}, \claimvalue{C63}] & Same test and fit counts \\
    Identity overlap $-$ joint, full pipeline & Matched support & \claimvalue{C50} &
      [\claimvalue{C51}, \claimvalue{C52}] & Includes prior and residual \\
    Max subject removal $-$ zero & Stimulus & \claimvalue{C55} &
      [\claimvalue{C56}, \claimvalue{C57}] & LaBraM intervention \\
    Subject removal $-$ equal-rank PCA & Stimulus & \claimvalue{C66} &
      [\claimvalue{C67}, \claimvalue{C68}] & Direction specificity \\
    Subject removal $-$ equal-rank random & Stimulus & \claimvalue{C69} &
      [\claimvalue{C70}, \claimvalue{C71}] & Direction specificity \\
    Condition prior $-$ EEG-only & External subject & \claimvalue{C27} &
      [\claimvalue{C28}, \claimvalue{C29}] & New subjects, repeated conditions \\
    \bottomrule
  \end{tabularx}
  % evidence: paper/generated/table_representation_comparisons.csv; paper/generated/table_external_validation.csv; paper/generated/table_reviewer_controls.csv
\end{table}

Under the same support-matched contrast, the semantic stimulus prior changes by
\claimvalue{C64} CCC and the hybrid stimulus prior by \claimvalue{C65}. These
branches are reported separately so the full-pipeline effect is not attributed
to EEG alone. Component ablations further separate nominal category, context,
source intensity code, and description text (Table~\ref{tab:prior-ablation}).

\begin{table}[t]
  \caption{Semantic-prior component audit. Positive values favor the full prior;
  intervals use the holdout identity required by each protocol.}
  \label{tab:prior-ablation}
  \centering
  \small
  \begin{tabularx}{\linewidth}{@{}llrr@{}}
    \toprule
    Protocol & Omitted component & Full $-$ ablated CCC & 95\% interval \cr
    \midrule
    Stimulus & Nominal category & \claimvalue{C87} & [\claimvalue{C88}, \claimvalue{C89}] \cr
    Subject & Description text & \claimvalue{C90} & [\claimvalue{C91}, \claimvalue{C92}] \cr
    Joint & Nominal category & \claimvalue{C93} & [\claimvalue{C94}, \claimvalue{C95}] \cr
    Joint & Context & \claimvalue{C96} & [\claimvalue{C97}, \claimvalue{C98}] \cr
    Joint & Intensity code & \claimvalue{C99} & [\claimvalue{C100}, \claimvalue{C101}] \cr
    Joint & Description text & \claimvalue{C102} & [\claimvalue{C103}, \claimvalue{C104}] \cr
    \bottomrule
  \end{tabularx}
  % evidence: paper/generated/table_reviewer_controls.csv
\end{table}

The strong nominal-category effect under unseen-stimulus and joint holdout
shows that the semantic prior's transferable score is largely experiment-label
structure. Description text contributes under subject holdout, whereas context,
intensity, and description do not have stable independent joint-holdout effects.
These are predictive ablations, not causal interventions on stimulus content.

\begin{table}[t]
  \caption{Identity-probe summary. Permuted values use the same representation,
  split, and estimator with shuffled training labels.}
  \label{tab:identity-probes}
  \centering
  \small
  \setlength{\tabcolsep}{4pt}
  \begin{tabularx}{\linewidth}{@{}Xlrrr@{}}
    \toprule
    Representation and task & Holdout & Observed & Permuted & Chance \\
    \midrule
    LaBraM pretrained, subject & Stimulus & \claimvalue{C09} & 0.0192 & \claimvalue{C46} \\
    LaBraM random, subject & Stimulus & \claimvalue{C10} & 0.0425 & \claimvalue{C46} \\
    DE, dataset & Subject & \claimvalue{C11} & 0.3913 & \claimvalue{C48} \\
    LaBraM pretrained, stimulus & Subject & \claimvalue{C12} & 0.0260 & \claimvalue{C49} \\
    External channel, condition & Subject & \claimvalue{C23} & \claimvalue{C24} & \claimvalue{C25} \\
    External global, condition & Subject & \claimvalue{C26} & 0.1692 & \claimvalue{C25} \\
    \bottomrule
  \end{tabularx}
  % evidence: paper/generated/table_identity_probes.csv; paper/generated/table_external_validation.csv
\end{table}

\section{External Validation Details}
\label{app:external-validation}

The initial subject-holdout analysis found that the condition prior exceeded
EEG-only macro-CCC by \claimvalue{C27}, with interval [\claimvalue{C28},
\claimvalue{C29}]. The v2 audit extends this case to exactly matched bandpower,
pretrained-LaBraM, and fixed-random-LaBraM representations under five protocols.
Table~\ref{tab:ds006866-v2-affect} reports every frozen cell rather than only
the favorable repeated-participant setting.

\begin{table*}[t]
  \caption{Complete ds006866 v2 affect audit. Entries are five-seed mean
  macro-CCC. BP denotes 320-dimensional channel bandpower; Pre. and Rand.
  denote frozen pretrained and fixed-random LaBraM. Hybrid is the transferable
  condition-family prior plus the EEG residual.}
  \label{tab:ds006866-v2-affect}
  \centering
  \scriptsize
  \setlength{\tabcolsep}{3.2pt}
  \begin{tabularx}{\textwidth}{@{}Xrrrrrrr@{}}
    \toprule
    Protocol & Prior & EEG BP & EEG Pre. & EEG Rand. & Hybrid BP & Hybrid Pre. & Hybrid Rand. \\
    \midrule
    Unseen participant & 0.3776 & 0.0369 & 0.0335 & 0.0198 & 0.3425 & 0.3742 & 0.3797 \\
    Nonsocial $\rightarrow$ social, repeated participant & 0.3757 & 0.2319 & 0.2604 & 0.2539 & 0.5144 & 0.5507 & 0.5476 \\
    Social $\rightarrow$ nonsocial, repeated participant & 0.3780 & 0.2226 & 0.2526 & 0.2455 & 0.5116 & 0.5486 & 0.5462 \\
    Nonsocial $\rightarrow$ social, unseen participant & 0.3704 & 0.0208 & 0.0323 & 0.0152 & 0.3318 & 0.3718 & 0.3726 \\
    Social $\rightarrow$ nonsocial, unseen participant & 0.3792 & 0.0172 & 0.0289 & 0.0220 & 0.3472 & 0.3761 & 0.3844 \\
    \bottomrule
  \end{tabularx}
  % evidence: paper/generated/ds006866_external_audit_v2.json
\end{table*}

For repeated participants, the hybrid improves over the prior in both
directions and all five participant-clustered MAE intervals lie below zero for
each representation. Under joint participant--context holdout, no LaBraM
hybrid interval establishes an improvement over the prior. Pretrained-minus-
random hybrid CCC is $-0.0008$ and $-0.0082$ across the two joint directions,
with only two of five seeds positive in each. Thus the strict result does not
support a pretraining-specific residual advantage.
% evidence: paper/generated/ds006866_external_audit_v2.json

\begin{table*}[t]
  \caption{Matched ds006866 v2 identity probes. Each cell is observed/permuted
  five-seed mean balanced accuracy. Chance reflects the fitted class set.}
  \label{tab:ds006866-v2-probes}
  \centering
  \scriptsize
  \setlength{\tabcolsep}{4pt}
  \begin{tabularx}{\textwidth}{@{}XXrrrr@{}}
    \toprule
    Protocol & Probe & BP & Pretrained & Random & Chance \\
    \midrule
    Unseen participant & Six-condition & 0.2448/0.1669 & 0.1713/0.1667 & 0.1675/0.1653 & 0.1667 \\
    Nonsocial $\rightarrow$ social, repeated participant & Participant & 0.9910/0.0066 & 0.9997/0.0074 & 0.9963/0.0057 & 0.0068 \\
    Social $\rightarrow$ nonsocial, repeated participant & Participant & 0.9890/0.0063 & 0.9999/0.0045 & 0.9973/0.0046 & 0.0068 \\
    Nonsocial $\rightarrow$ social, unseen participant & Condition family & 0.4091/0.3238 & 0.3399/0.3317 & 0.3306/0.3333 & 0.3333 \\
    Social $\rightarrow$ nonsocial, unseen participant & Condition family & 0.4206/0.3268 & 0.3384/0.3336 & 0.3308/0.3315 & 0.3333 \\
    \bottomrule
  \end{tabularx}
  % evidence: paper/generated/ds006866_external_audit_v2.json
\end{table*}

The near-perfect repeated-participant probes establish information available
to a fixed linear readout; they do not show that the affect head causally uses
identity. Public events expose condition but not trial-level image identity, so
stimulus holdout remains blocked. Header-only conformance confirms the same
ordered 64-channel EEG set in all 148 recordings; 146 match the base source
pattern and two contain one additional, predeclared EMG channel.
% evidence: paper/generated/ds006866_channel_audit_v2.json

\section{Authorized SEED-V Protocol-Audit Case Study}
\label{app:seedv}

SEED-V is a stimulus-induced five-class EEG and eye-movement study with
16 participants, three sessions, and fifteen film clips per session
\citep{li2019seedv}. We report a separately authorized raw-EEG analysis as an
independent case study rather than a fifth OpenAffect-EEG source. Its source
recordings, feature arrays, media, trial predictions, participant codes, and
project paths are not redistributed. The public artifact contains only the
source-safe aggregate in \texttt{paper/generated/seedv\_case\_study.json},
its input and code hashes, and this interpretation boundary.

The authorized project compared a raw-temporal baseline against a pre-existing
SDSA variant under a strict leave-one-participant-out protocol and a joint
participant--stimulus protocol. Table~\ref{tab:seedv-case} reports all frozen
cells rather than selecting the favorable joint cell. The strict SDSA contrast
is negative, and both bootstrap intervals include zero. Its frozen decision
gate therefore fails; the contrast is evidence of protocol sensitivity, not a
model-improvement claim. Separately, the authorized project's strict-to-
PaperStyle ladder rises from 0.2786 to 0.5922, so the large protocol gap cannot
be attributed to a model variant alone \citep{li2019seedv}.

\begin{table}[t]
  \caption{Authorized SEED-V aggregate audit. Values are mean trial macro-F1
  across the frozen result rows; $\Delta$ is SDSA minus raw temporal. No
  participant-level or trial-level values are released.}
  \label{tab:seedv-case}
  \centering
  \small
  \setlength{\tabcolsep}{5pt}
  \begin{tabularx}{\linewidth}{@{}Xrrrr@{}}
    \toprule
    Protocol & Raw temporal & SDSA & $\Delta$ & 95\% interval \\
    \midrule
    Strict participant holdout & 0.2418 & 0.2297 & -0.0121 & [-0.0439, 0.0178] \\
    Joint participant--stimulus holdout & 0.1788 & 0.1856 & 0.0068 & [-0.0154, 0.0292] \\
    \bottomrule
  \end{tabularx}
  % evidence: paper/generated/seedv_case_study.json
\end{table}

\section{Independent MusicEEG Protocol Audits}
\label{app:music-labram}

To test whether the EmoEEG-MC representation finding is confined to one visual
paradigm, we run two prespecified audits on independent MusicEEG
\citep{daly2020musiceeg}. Each uses 1,240 labelled 12-second music trials from
31 participants, five frozen complete-trial splits, and source-safe aggregate
reporting. First, a transparent channel log-bandpower Ridge readout reaches
macro CCC 0.1031 under random-trial evaluation, 0.0175 with unseen
participants, 0.1130 with unseen music but known participants, and 0.0227
under joint participant--music holdout. The paired random-minus-joint
difference is 0.0804 (seed SD 0.0357). Thus holding out music alone does not
substitute for holding out participants in this auditory source.

The same bandpower representation identifies training-seen participants under
unseen music at balanced accuracy 0.9995 (permuted 0.0431; 31-class uniform
chance 0.0323). In contrast, the seen-class music probe under unseen
participants is 0.0000 balanced accuracy (permuted 0.0033; 78.8--91.3\% test
coverage across seeds). These probes establish available information under
their stated coverage, not causal use by the affect predictor.
% evidence: paper/generated/music_eeg_external_audit_v2.json

The self-report diagnostic also constrains interpretation. Leave-one-rating-out
same-music consensus CCC is 0.1437 for valence and 0.1634 for arousal (1,123
supported ratings each). Descriptive in-sample additive models assign more
R-squared to music identity (0.3579 and 0.3484) than participant identity
(0.1170 and 0.1782); they are neither test--retest reliability estimates nor
causal variance decompositions.
% evidence: paper/generated/music_eeg_external_audit_v2.json

Second, we apply the frozen official LaBraM base encoder to exactly the same
MusicEEG split family. Each trial uses its 19 recorded 10--20 channels and its
observed event window; MNE volts are converted to microvolts, common-average
referenced within trial, resampled to 200~Hz, and final tokens pooled. There is
no padding, cross-trial window, fine-tuning, or feature release. Pretrained and
fixed-random initializations use the same 200-dimensional Ridge readout,
manifests, five split seeds, and nested validation selection.
% constants-source: cards/datasets/ds002721.md; scripts/extract_labram_bids_features.py

Table~\ref{tab:music-labram} is a separate foundation-model audit rather
than a new model claim. Pretraining has a positive matched macro-CCC point difference
in all four audited protocols, including participant and joint holdout, but the
corresponding unseen-music subject probe is also substantially higher for
pretrained features. This external source therefore qualifies, rather than
reverses, the main case study: pretrained features have higher report-prediction
point estimates in this auditory source while retaining participant structure, so neither
the positive gain nor a single identity probe supports a general invariant-
emotion conclusion.

\begin{table}[t]
  \caption{Independent MusicEEG frozen LaBraM audit. Values summarize five
  fixed split seeds; $\Delta$ is pretrained minus fixed-random macro-CCC.
  Subject probes predict training-seen participants under unseen music.}
  \label{tab:music-labram}
  \centering
  \small
  \setlength{\tabcolsep}{4pt}
  \begin{tabularx}{\linewidth}{@{}Xrrr@{}}
    \toprule
    Protocol & Pretrained CCC & Random CCC & $\Delta$ CCC \\
    \midrule
    Random trial & 0.0923 & 0.0191 & 0.0732 \\
    Unseen participant & 0.0697 & -0.0006 & 0.0703 \\
    Unseen music & 0.0847 & 0.0145 & 0.0703 \\
    Unseen participant and music & 0.0422 & 0.0085 & 0.0337 \\
    \midrule
    Unseen-music subject probe (BA) & 0.9136 & 0.4239 & --- \\
    Permuted subject probe (BA) & 0.0313 & 0.0297 & chance 0.0323 \\
    \bottomrule
  \end{tabularx}
  % evidence: paper/generated/music_eeg_external_audit_v2.json
\end{table}

\section{Compute and Software Environment}

Transparent features, linear baselines, probes, paired statistics, artifact
compilation, and tests run in a CPU environment on Linux x86-64 with Python
3.11, NumPy 2.4.6, pandas 3.0.5, SciPy 1.17.1, scikit-learn 1.9.0, and MNE
1.12.1. Foundation feature extraction and EEGNet use one NVIDIA RTX 6000 Ada
Generation GPU with 49,140 MiB memory, PyTorch 2.11.0, and CUDA 13.0. CPU-only
adversarial residual runs use four threads to avoid competing with GPU jobs.
% evidence: docs/reproduction_report.md; configs/recompute_neurips_ed_v1.yaml

The typed recomputation graph separates feature extraction, evaluation,
statistics, and manuscript compilation. Exact historical wall-clock and energy
totals were not captured consistently for every preliminary or failed run; the
release therefore reports component environments and DAG boundaries rather
than inventing a project-wide total. The final cold-run release log is designed
to record per-step wall times and will be published with the executable
artifact. This limitation does not affect frozen scores but constrains precise
compute-cost replication.

\section{Release, Licensing, and Reproduction Boundaries}

The anonymous package contains code, tests, configurations, evaluation cards,
dataset and model cards, Croissant metadata, generated table sources, and claim
configuration. It explicitly excludes raw EEG, source feature arrays,
foundation checkpoints, movie, music, and image stimuli, office documents, Git
history, and machine-local paths.

\texttt{metadata} mode validates the benchmark contract without source data.
\texttt{smoke} uses deterministic synthetic trials to exercise splitting,
baselines, probes, and a small registry. \texttt{full} requires an explicit
user-provided data root and verifies frozen input hashes and paper assets. The
typed \texttt{recompute plan/run} interface covers trial construction, feature
extraction, evaluation, probes, statistics, and paper compilation;
\texttt{--resume} skips a step only when all declared outputs exist.

The shortest exact verification sequence from a clean checkout is:
\begin{small}
\begin{verbatim}
conda env create -f environment.yml
conda activate openaffect-eeg
python -m pip install -e ".[smoke]"
python scripts/run_benchmark.py metadata \
  --output-root /tmp/openaffect-metadata
python scripts/run_benchmark.py smoke \
  --output-root /tmp/openaffect-smoke
export D=/path/to/OpenAffect-EEG-data
python scripts/run_benchmark.py full \
  --data-root "$D" --output-root "$D/reproduction"
python scripts/build_neurips_submission.py
\end{verbatim}
\end{small}
The first two modes never read EEG. Full mode verifies the frozen registry and
rebuilds evidence-linked tables; it is not mislabeled as raw feature extraction
or model retraining. The separately documented \texttt{recompute\_benchmark.py
plan/run} commands expose that typed raw-to-paper DAG and require explicit paths
to source data, checkpoints, and foundation-model repositories.

Full-mode verification of frozen outputs is not described as a cold download
and multi-day retraining run. The existing isolated-environment verification
was performed on the same compute host and is not represented as independent
laboratory reproduction.

\section{Machine-Readable Attachments}

The full registry retains per-seed identity results, all paired model
comparisons, band and region sensitivity cells, prediction-key audits, source
hashes, and unsupported protocol states. These files prevent selective
compression of null or unfavorable results into the PDF. Region analyses are
predictive sensitivity tests and are not interpreted as neural source
localization.

\section{Position Relative to Closest Evaluation Resources}

\begin{table}[t]
  \caption{The closest resources address complementary layers of the problem.
  The final column states the additional evaluation object introduced here,
  rather than asserting that prior work is invalid.}
  \label{tab:related-benchmarks}
  \centering
  \small
  \begin{tabularx}{\linewidth}{@{}lXX@{}}
    \toprule
    Resource & Primary evaluation object & Additional OpenAffect-EEG object \\
    \midrule
    Leakage audit \citep{lei2025leakage} & Trial-segment and test-data leakage on DEAP & Subject, stimulus, joint, context, and dataset identity contracts \\
    mdJPT \citep{zhang2025mdjpt} & Multi-dataset affect pretraining and new-dataset transfer & Deployment-specific identity isolation plus matched random encoders \\
    EMOD \citep{chen2026emod} & Discrete/continuous label alignment in a common V--A space & Self-report prior/residual separation and stimulus provenance states \\
    EEG-FM-Compass \citep{liu2026compass} & Cross-task FM comparison under subject transfer and calibration & Affect-specific stimulus and condition audits with predictor-linked interventions \\
    FMScope \citep{lin2026identitytrap} & Subject-axis diagnosis and erasure across clinical EEG-FMs & Joint subject--stimulus--context contracts and report-level affect estimands \\
    EEG concept audit \citep{tang2026capture} & Encoded-versus-used physiological concepts under erasure & Experiment-metadata priors separated from trial-level self-reports \\
    Emotion protocol audit \citep{suo2026protocols} & Subject-dependent versus subject-disjoint SEED evaluation & Stimulus, joint, context, and dataset isolation with provenance states \\
    OpenAffect-EEG & Executable evidence and applicability contract & Hash registry, claim ledger, cards, blocked cells, and matched controls \\
    \bottomrule
  \end{tabularx}
\end{table}
\placeholderfigure{figures/figureS_condition_audit.pdf}{2.00in}{Initial
ds006866 subject-holdout condition audit. The condition prior, direct EEG, and
condition-plus-EEG residual are reported separately; paired residual effects and
condition probes use the five fixed seeds. Image identity is unavailable, so
this is condition-level evidence and cannot support a stimulus-holdout claim.}
{fig:condition-audit}
